% ============================================================
%  FGBench
% ============================================================
\section{FGBench: Benchmark Construction}
\label{sec:fgbench}

FGBench asks which of two unrelated pairs is noisier, and its construction follows
from that question. Every image and caption is an existing COCO~\cite{lin2014coco} or
Visual Genome~\cite{krishna2017vg} item; a mismatch is produced by contradicting a
structured fact attached to the image, so the label rests on an external
justification. Splits are disjoint at the cycle level, and the construction is
auditable.

For an image pair $(A,B)$ with captions $(T_A,T_B)$ we emit four records
\begin{equation}
\begin{aligned}
  &(A,T_A,\text{clean}),\quad (B,T_A,\text{noisy}),\\
  &(B,T_B,\text{clean}),\quad (A,T_B,\text{noisy}),
\end{aligned}
\end{equation}
a reciprocal cycle. Each image and caption occurs once as a match and once as a
mismatch, so every facet is exactly $50\%$ balanced and the two classes share the
same distributions over image and text sides. A detector cannot exploit the frequency
of an image or caption to separate the classes. The cycle is a property of the data,
not of the method: CRD is trained on individual pairs and scores one pair at
inference.

Five facets come from two sources. From COCO we build \textbf{Entity} (object
category substituted) and \textbf{Quantity} (numeral substituted); from Visual Genome
we build \textbf{Attribute} (colour, material or state), \textbf{Action} (predicate)
and \textbf{Relation} (spatial or comparative relation), all from VG150-style
annotations.

\begin{table}[t]
\centering
\small
\setlength{\tabcolsep}{4pt}
\caption{FGBench v1-strict. Every cycle contributes two matched and two mismatched
pairs, so each facet is exactly balanced. The test split, used by all methods,
contains $17{,}565$ cycles and $70{,}260$ pairs; per-facet test sizes make the
statistical power of each column explicit.}
\label{tab:fgbench}
\begin{tabular}{llrrr}
\toprule
Source & Facet & Cycles & Pairs & Test (cycles) \\
\midrule
COCO          & Entity     & 32{,}589 & 130{,}356 &  2{,}298 \\
COCO          & Quantity   &  2{,}940 &  11{,}760 &    114 \\
Visual Genome & Attribute  & 49{,}042 & 196{,}168 & 14{,}281 \\
Visual Genome & Action     &    567   &   2{,}268 &    166 \\
Visual Genome & Relation   &  2{,}886 &  11{,}544 &    706 \\
\midrule
\multicolumn{2}{l}{Total}      & 88{,}024 & 352{,}096 & 17{,}565 \\
\bottomrule
\end{tabular}
\end{table}

Table~\ref{tab:fgbench} summarizes the construction. Labels are \emph{silver}: they
follow from a substitution rule applied to a structured annotation, not from human
adjudication. An annotation records what was
annotated, not a complete inventory of the image, so unannotated content is treated
as unknown rather than absent; a strict tier retains only Relation cases with
unambiguous annotation support and removes $276$ cycles ($1{,}104$ rows) whose labels
rely on six ambiguous rules; all numbers here refer to v1-strict. All detectors are
evaluated on the identical manifest. The benchmark audit reports a maximum of one
split per image and per caption and evidence completeness $1.000$. Because every
image and caption occurs once in each class, any single-side score has identical
class-conditional distributions; the TF-IDF check is therefore an analytic invariant
($0.5000$), not independent evidence, and the construction makes a text-only shortcut
impossible by design.
